% Part: first-order-logic
% Chapter: natural-deduction
% Section: soundness

\documentclass[../../../include/open-logic-section]{subfiles}

\begin{document}

\iftag{FOL}
      {\olfileid{fol}{ntd}{sou}}
      {\olfileid{pl}{ntd}{sou}}
      
\olsection{Correctheid}

\begin{explain}
!!^a{derivation}ssysteem, zoals natuurlijke deductie, is
\emph{correct} als het geen zaken kan !!{derive} die in werkelijkheid
niet volgen. Correctheid is dus een soort gegarandeerde
veiligheidseigenschap van !!{derivation}ssystemen. Afhankelijk van de
bewijstheoretische eigenschap in kwestie willen we bijvoorbeeld weten
dat
\begin{enumerate}
\item elke !!{derivable} !!{sentence} \iftag{FOL}{geldig}{een
  tautologie} is;
\item als !!a{sentence} !!{derivable} is uit enkele andere, zij ook een
  gevolg daarvan is;
\item als een verzameling !!{sentence}s inconsistent is, zij
  onvervulbaar is.
\end{enumerate}
Dit zijn belangrijke eigenschappen van !!a{derivation}ssysteem. Als
een ervan niet geldt, schiet het !!{derivation}ssysteem tekort: het zou
te veel !!{derive}. Het is daarom van het grootste belang de
correctheid van !!a{derivation}ssysteem vast te stellen.
\end{explain}

\begin{thm}[Correctheid]
\ollabel{thm:soundness}
Als $!A$ !!{derivable} is uit de !!{undischarged} aannamen
$\Gamma$, dan $\Gamma \Entails !A$.
\end{thm}

\begin{proof}
Laat $\delta$ !!a{derivation} van $!A$ zijn. We voeren inductie uit
naar het aantal inferenties in~$\delta$.

Voor de inductiebasis tonen we de bewering aan wanneer het aantal
inferenties~$0$ is. In dat geval bestaat $\delta$ uitsluitend uit één
!!{sentence}~$!A$, dus uit een aanname. Die aanname is
!!{undischarged}, want aannamen kunnen alleen door inferenties worden
!!{discharged}, en er zijn geen inferenties. Elke
\iftag{FOL}{!!{structure}~$\Struct{M}$}{!!{valuation}~$\pAssign{v}$}
die alle !!{undischarged} aannamen van het bewijs waar maakt, maakt
daarom ook~$!A$ waar.

Nu volgt de inductiestap. Veronderstel dat $\delta$~$n$ inferenties
bevat. De premisse of premissen van de onderste inferentie worden
!!{derive}d met deel-!!{derivation}s die elk minder dan~$n$
inferenties bevatten. We nemen als inductiehypothese aan dat de
premissen van de onderste inferentie volgen uit de !!{undischarged}
aannamen van de deel-!!{derivation}s die in die premissen eindigen. We
moeten aantonen dat de conclusie~$!A$ volgt uit de !!{undischarged}
aannamen van het gehele bewijs.

We onderscheiden gevallen naar het soort onderste inferentie. Eerst
beschouwen we de mogelijke inferenties met slechts één premisse.

\begin{enumerate}
\item Veronderstel dat de laatste inferentie \Intro{\lnot} is. De
  !!{derivation} heeft dan de vorm
  \begin{prooftree}
    \AxiomC{$\Gamma, \Discharge{!A}{n}$}
    \RightLabel{$\delta_1$}
    \DeduceC{$\lfalse$}
    \DischargeRule{\Intro{\lnot}}{n}
    \UnaryInfC{$\lnot !A$}
  \end{prooftree}
  Volgens de inductiehypothese volgt $\lfalse$ uit de
  !!{undischarged} aannamen $\Gamma \cup \{!A\}$ van~$\delta_1$.
  Beschouw
\iftag{FOL}{!!a{structure}~$\Struct{M}$}{!!a{valuation}~$\pAssign{v}$}.
  We moeten aantonen dat, als
  $\iftag{FOL}{\Sat{M}}{\pSat{v}}{\Gamma}$, dan
  $\iftag{FOL}{\Sat{M}}{\pSat{v}}{\lnot !A}$. Veronderstel voor een
  bewijs uit het ongerijmde dat
  $\iftag{FOL}{\Sat{M}}{\pSat{v}}{\Gamma}$, maar
  $\iftag{FOL}{\Sat/{M}}{\pSat/{v}}{\lnot !A}$, dat wil zeggen
  $\iftag{FOL}{\Sat{M}}{\pSat{v}}{!A}$. Dit zou betekenen dat
  $\iftag{FOL}{\Sat{M}}{\pSat{v}}{\Gamma \cup \{!A\}}$, in strijd
  met onze inductiehypothese. Dus
  $\iftag{FOL}{\Sat{M}}{\pSat{v}}{\lnot !A}$.

\item De laatste inferentie is \Elim{\land}. Er zijn twee varianten:
  $!A$ of $!B$ mag uit de premisse $!A \land !B$ worden afgeleid.
  Beschouw het eerste geval. De !!{derivation}~$\delta$ ziet er als
  volgt uit:
  \begin{prooftree}
    \AxiomC{$\Gamma$}
    \RightLabel{$\delta_1$}
    \DeduceC{$!A \land !B$}
    \RightLabel{\Elim{\land}}
    \UnaryInfC{$!A$}
  \end{prooftree}
  Volgens de inductiehypothese volgt $!A \land !B$ uit de
  !!{undischarged} aannamen~$\Gamma$ van~$\delta_1$. Beschouw
  !!a{structure}~\iftag{FOL}{$\Struct{M}$}{$\pAssign{v}$}. We moeten
  aantonen dat, als $\iftag{FOL}{\Sat{M}}{\pSat{v}}{\Gamma}$, dan
  $\iftag{FOL}{\Sat{M}}{\pSat{v}}{!A}$. Veronderstel
  $\iftag{FOL}{\Sat{M}}{\pSat{v}}{\Gamma}$. Uit onze
  inductiehypothese ($\Gamma \Entails !A \land !B$) weten we dat
  $\iftag{FOL}{\Sat{M}}{\pSat{v}}{!A \land !B}$. Per definitie geldt
  $\iftag{FOL}{\Sat{M}}{\pSat{v}}{!A \land !B}$ dan en slechts dan
  als $\iftag{FOL}{\Sat{M}}{\pSat{v}}{!A}$ en
  $\iftag{FOL}{\Sat{M}}{\pSat{v}}{!B}$. (Het geval waarin $!B$ uit
  $!A \land !B$ wordt afgeleid, verloopt op dezelfde wijze.)
  
\item De laatste inferentie is \Intro{\lor}. Er zijn twee varianten:
  $!A \lor !B$ mag worden afgeleid uit de premisse~$!A$ of uit de
  premisse~$!B$. Beschouw het eerste geval. De !!{derivation} heeft de
  vorm
  \begin{prooftree}
    \AxiomC{$\Gamma$}
    \RightLabel{$\delta_1$}
    \DeduceC{$!A$}
    \RightLabel{\Intro{\lor}}
    \UnaryInfC{$!A \lor !B$}
  \end{prooftree}
  Volgens de inductiehypothese volgt $!A$ uit de !!{undischarged}
  aannamen~$\Gamma$ van~$\delta_1$. Beschouw
\iftag{FOL}{!!a{structure}~$\Struct{M}$}{!!a{valuation}~$\pAssign{v}$}.
  We moeten aantonen dat, als
  $\iftag{FOL}{\Sat{M}}{\pSat{v}}{\Gamma}$, dan
  $\iftag{FOL}{\Sat{M}}{\pSat{v}}{!A \lor !B}$. Veronderstel
  $\iftag{FOL}{\Sat{M}}{\pSat{v}}{\Gamma}$. Dan
  $\iftag{FOL}{\Sat{M}}{\pSat{v}}{!A}$, omdat
  $\Gamma \Entails !A$ (de inductiehypothese). Bijgevolg geldt ook
  $\iftag{FOL}{\Sat{M}}{\pSat{v}}{!A \lor !B}$. (Het geval waarin
  $!A \lor !B$ uit~$!B$ wordt afgeleid, verloopt op dezelfde wijze.)
  
\item De laatste inferentie is \Intro{\lif}. $!A \lif !B$ wordt
  afgeleid uit een deelbewijs met aanname~$!A$ en conclusie~$!B$, dat
  wil zeggen:
  \begin{prooftree}
    \AxiomC{$\Gamma, \Discharge{!A}{n}$}
    \RightLabel{$\delta_1$}
    \DeduceC{$!B$}
    \DischargeRule{\Intro{\lif}}{n}
    \UnaryInfC{$!A \lif !B$}
  \end{prooftree}
  Volgens de inductiehypothese volgt $!B$ uit de !!{undischarged}
  aannamen van~$\delta_1$, dat wil zeggen:
  $\Gamma \cup \{!A\} \Entails !B$. Beschouw
\iftag{FOL}{!!a{structure}~$\Struct{M}$}{!!a{valuation}~$\pAssign{v}$}.
  De !!{undischarged} aannamen van~$\delta$ zijn precies $\Gamma$,
  want $!A$ wordt bij de laatste inferentie opgeheven. We moeten dus
  aantonen dat $\Gamma \Entails !A \lif !B$. Veronderstel voor een
  bewijs uit het ongerijmde dat er een
  \iftag{FOL}{!!{structure}~$\Struct{M}$}{!!{valuation}~$\pAssign{v}$},
  is waarvoor $\iftag{FOL}{\Sat{M}}{\pSat{v}}{\Gamma}$ maar
  $\iftag{FOL}{\Sat/{M}}{\pSat/{v}}{!A \lif !B}$. Dan
  $\iftag{FOL}{\Sat{M}}{\pSat{v}}{!A}$ en
  $\iftag{FOL}{\Sat/{M}}{\pSat/{v}}{!B}$. Maar volgens de hypothese is
  $!B$ een gevolg van $\Gamma \cup \{!A\}$, dat wil zeggen:
  $\iftag{FOL}{\Sat{M}}{\pSat{v}}{!B}$, en dat is een tegenspraak.
  Dus
  $\Gamma \Entails !A \lif !B$.

\item De laatste inferentie is \FalseInt. Hier eindigt $\delta$ op
  \begin{prooftree}
    \AxiomC{$\Gamma$}
    \RightLabel{$\delta_1$}
    \DeduceC{$\lfalse$}
    \RightLabel{\FalseInt}
    \UnaryInfC{$!A$}
  \end{prooftree}
  Volgens de inductiehypothese is $\Gamma \Entails \lfalse$. We
  moeten aantonen dat $\Gamma \Entails !A$. Veronderstel dat dit niet
  zo is. Dan geldt voor een zekere
  ~$\iftag{FOL}{\Struct{M}}{\pAssign{v}}$:
    $\iftag{FOL}{\Sat{M}}{\pSat{v}}{\Gamma}$ en
    $\iftag{FOL}{\Sat/{M}}{\pSat/{v}}{!A}$. We hebben echter altijd
    $\iftag{FOL}{\Sat/{M}}{\pSat/{v}}{\lfalse}$. Dit zou dus
    $\Gamma \Entails/ \lfalse$ betekenen, in strijd met de
    inductiehypothese.

\item De laatste inferentie is \FalseCl. Oefening.

\iftag{FOL}{
\item De laatste inferentie is \Intro{\lforall}. Dan heeft $\delta$ de
  vorm
  \begin{prooftree}
    \AxiomC{$\Gamma$}
    \RightLabel{$\delta_1$}
    \DeduceC{$!A(a)$}
    \RightLabel{\Intro{\lforall}}
    \UnaryInfC{$\lforall[x][!A(x)]$}
  \end{prooftree}
  Volgens de inductiehypothese is de premisse $!A(a)$ een gevolg van
  de !!{undischarged} aannamen $\Gamma$. Beschouw een structuur
  $\Struct{M}$ waarvoor $\Sat{M}{\Gamma}$. We moeten aantonen dat
  $\Sat{M}{\lforall[x][!A(x)]}$. Omdat $\lforall[x][!A(x)]$
  !!a{sentence} is, moeten we dus aantonen dat voor elke
  variabelentoekenning~$s$ geldt: $\Sat{M}{!A(x)}[s]$
  (\olref[syn][ass]{prop:sat-quant}). Omdat $\Gamma$ uitsluitend uit
  zinnen bestaat, geldt $\Sat{M}{!B}[s]$ voor elke
  $!B \in \Gamma$, volgens
  \olref[syn][sat]{defn:satisfaction}. Laat $\Struct{M'}$ gelijk zijn
  aan $\Struct{M}$,
  behalve dat $\Assign{a}{M'} = s(x)$. Omdat $a$ niet in~$\Gamma$
  voorkomt, geldt volgens
  \olref[syn][ext]{cor:extensionality-sent} dat
  $\Sat{M'}{\Gamma}$. Omdat $\Gamma \Entails !A(a)$, geldt
  $\Sat{M'}{!A(a)}$. Omdat $!A(a)$ !!a{sentence} is, geldt
  $\Sat{M'}{!A(a)}[s]$ volgens
  \olref[syn][ass]{prop:sentence-sat-true}.
  $\Sat{M'}{!A(x)}[s]$ geldt dan en slechts dan als
  $\Sat{M'}{!A(a)}$, volgens
  \olref[syn][ext]{prop:ext-formulas} (bedenk dat $!A(a)$ niets
  anders is dan $\Subst{!A(x)}{a}{x}$). Dus
  $\Sat{M'}{!A(x)}[s]$. Omdat $a$ niet
  in~$!A(x)$ voorkomt, volgt uit
  \olref[syn][ext]{prop:extensionality} dat
  $\Sat{M}{!A(x)}[s]$. Maar $s$ was een willekeurige
  variabelentoekenning, en dus $\Sat{M}{\lforall[x][!A(x)]}$.
  
\item De laatste inferentie is \Intro{\lexists}. Oefening.

\item De laatste inferentie is \Elim{\forall}. Oefening.
}{}
\end{enumerate}

Laten we nu de mogelijke inferenties met meerdere premissen
beschouwen: \Elim{\lor}, \Intro{\land},
\iftag{FOL}{\Elim{\lif} en \Elim{\lexists}}{en \Elim{\lif}}.
\begin{enumerate}

\item De laatste inferentie is \Intro{\land}. $!A \land !B$ wordt
  afgeleid uit de premissen $!A$ en $!B$, en $\delta$ heeft de vorm
  \begin{prooftree}
    \AxiomC{$\Gamma_1$}
    \RightLabel{$\delta_1$}
    \DeduceC{$!A$}
    \AxiomC{$\Gamma_2$}
    \RightLabel{$\delta_2$}
    \DeduceC{$!B$}
    \RightLabel{\Intro{\land}}
    \BinaryInfC{$!A \land !B$}
  \end{prooftree}
  Volgens de inductiehypothese volgt $!A$ uit de !!{undischarged}
  aannamen~$\Gamma_1$ van~$\delta_1$, en volgt $!B$ uit de
  !!{undischarged} aannamen~$\Gamma_2$ van~$\delta_2$. De
  !!{undischarged} aannamen van~$\delta$ zijn $\Gamma_1 \cup
  \Gamma_2$. We moeten dus aantonen dat
  $\Gamma_1 \cup \Gamma_2 \Entails !A \land !B$. Beschouw
  \iftag{FOL}{!!a{structure}~$\Struct{M}$}{!!a{valuation}~$\pAssign{v}$}
  waarvoor
  $\iftag{FOL}{\Sat{M}}{\pSat{v}}{\Gamma_1 \cup \Gamma_2}$. Omdat
  $\iftag{FOL}{\Sat{M}}{\pSat{v}}{\Gamma_1}$, moet
  $\iftag{FOL}{\Sat{M}}{\pSat{v}}{!A}$ gelden, aangezien
  $\Gamma_1 \Entails !A$. En omdat
  $\iftag{FOL}{\Sat{M}}{\pSat{v}}{\Gamma_2}$, geldt
  $\iftag{FOL}{\Sat{M}}{\pSat{v}}{!B}$, aangezien
  $\Gamma_2 \Entails !B$. Samen leveren deze feiten
  $\iftag{FOL}{\Sat{M}}{\pSat{v}}{!A \land !B}$ op.
  
\item De laatste inferentie is \Elim{\lor}. Oefening.

\item De laatste inferentie is \Elim{\lif}. $!B$ wordt afgeleid uit
  de premissen $!A \lif !B$ en~$!A$. De !!{derivation}~$\delta$ ziet
  er als volgt uit:
  \begin{prooftree}
    \AxiomC{$\Gamma_1$}
    \RightLabel{$\delta_1$}
    \DeduceC{$!A \lif !B$}
    \AxiomC{$\Gamma_2$}
    \RightLabel{$\delta_2$}
    \DeduceC{$!A$}
    \RightLabel{\Elim{\lif}}
    \BinaryInfC{$!B$}
  \end{prooftree}
  Volgens de inductiehypothese volgt $!A \lif !B$ uit de
  !!{undischarged} aannamen~$\Gamma_1$ van~$\delta_1$, en volgt $!A$
  uit de !!{undischarged} aannamen~$\Gamma_2$ van~$\delta_2$.
  Beschouw
\iftag{FOL}{!!a{structure}~$\Struct{M}$}{!!a{valuation}~$\pAssign{v}$}.
  We moeten aantonen dat, als
  $\iftag{FOL}{\Sat{M}}{\pSat{v}}{\Gamma_1 \cup \Gamma_2}$, dan
  $\iftag{FOL}{\Sat{M}}{\pSat{v}}{!B}$. Veronderstel
  $\iftag{FOL}{\Sat{M}}{\pSat{v}}{\Gamma_1 \cup \Gamma_2}$. Omdat
  $\Gamma_1 \Entails !A \lif !B$, geldt
  $\iftag{FOL}{\Sat{M}}{\pSat{v}}{!A \lif !B}$. Omdat
  $\Gamma_2 \Entails !A$, geldt
  $\iftag{FOL}{\Sat{M}}{\pSat{v}}{!A}$. Dit betekent dat
  $\iftag{FOL}{\Sat{M}}{\pSat{v}}{!B}$ (want als
  $\iftag{FOL}{\Sat/{M}}{\pSat/{v}}{!B}$, dan zou uit
  $\iftag{FOL}{\Sat{M}}{\pSat{v}}{!A}$ volgen dat
  $\iftag{FOL}{\Sat/{M}}{\pSat/{v}}{!A \lif !B}$,
  in tegenspraak met
  ~$\iftag{FOL}{\Sat{M}}{\pSat{v}}{!A \lif !B}$).

\item De laatste inferentie is \Elim{\lnot}. Oefening.
  
\tagitem{FOL}{De laatste inferentie is \Elim{\lexists}. Oefening.}{}
\end{enumerate}
\end{proof}

\tagprob{FOL}
\begin{prob}
Voltooi het bewijs van \olref[fol][ntd][sou]{thm:soundness}.
\end{prob}
\tagendprob

\tagprob{notFOL}
\begin{prob}
Voltooi het bewijs van \olref[pl][ntd][sou]{thm:soundness}.
\end{prob}
\tagendprob

\begin{cor}
\ollabel{cor:weak-soundness}
Als $\Proves !A$, dan is $!A$ \iftag{FOL}{geldig}{een tautologie}.
\end{cor}

\begin{cor}
\ollabel{cor:consistency-soundness}
Als $\Gamma$ vervulbaar is, dan is zij consistent.
\end{cor}

\begin{proof}
We bewijzen de contrapositie. Veronderstel dat $\Gamma$ niet
consistent is. Dan $\Gamma \Proves \lfalse$, dat wil zeggen: er bestaat
!!a{derivation} van $\lfalse$ uit !!{undischarged} aannamen
in~$\Gamma$. Volgens \olref{thm:soundness} moet elke
\iftag{FOL}{!!{structure}~$\Struct{M}$}{!!{valuation}~$\pAssign{v}$}
die $\Gamma$ waar maakt, ook $\lfalse$ waar maken. Omdat
$\iftag{FOL}{\Sat/{M}}{\pSat/{v}}{\lfalse}$ voor elke
\iftag{FOL}{!!{structure}~$\Struct{M}$}{!!{valuation}~$\pAssign{v}$},
kan geen enkele \iftag{FOL}{$\Struct{M}$}{$\pAssign{v}$} $\Gamma$
waar maken. $\Gamma$ is dus niet vervulbaar.
\end{proof}
\end{document}
